\section{A uniform zero representation and constructed majors}
\label{truthy:section}

The new contribution is a change of proof presentation.  The numerical
operations, native formulas, and primitive mathematical rules remain
those specified above.  In particular, this section neither removes an
arithmetical axiom nor changes the domain of interpretation.  It changes
which finite data represent zero and how an elimination packet specifies
the theorem that must be proved.

The distinction is useful here.  Earlier packets used Python falsiness
both as an instruction marker and as the canonical spelling of
arithmetical zero.  A zero initializer or predecessor consequently needed
a separate proof replacing it by a nonempty term.  We choose a canonical
zero that is already truthy, maintain that convention during opening,
and construct an equality or closure major from its component fields.
The construction still checks a proof of that major.

\subsection{Canonical expressions and the zero convention}

Retain the inherited arity codec.  Under \(p\) protected numerical slots,
formula depth is \(2p\), and numerical depth is \(3+2p\).
At a numerical position, coordinates \(0,1,2\) denote zero, while
\(3+2j,4+2j\) denote protected slot \(j\), counted from the inside.
Admitted negative integers are rigid parameters.  The meaningful
constructors and their environments are as follows:
\[
\begin{array}{c|c|l}
\text{sort}&\text{raw constructor}&\text{additional child environment}\\ \hline
\text{number}&[n]&\text{two dummy zero slots}\\
\text{number}&[n,a,h]&
 \text{one zero slot for }n,a;\ \text{one state slot for }h\\
\text{formula}&[s,t]&\text{the same numerical slots}\\
\text{formula}&[B]&\text{the two slots of }\D\\
\text{formula}&[B,B,B]&\text{the one slot of }\Q .
\end{array}
\]
Thus the common depth increment \(5-\operatorname{len}\) retains its
existing meaning.  Numerical binary and surplus arrays denote zero;
formula surplus arrays and admitted scalar formula padding denote the
fixed truth \(0=0\).  Those aliases are administrative syntax, not new
arithmetical or logical symbols.

In the new canonical codec, \emph{raw scalar \(1\) denotes numerical
zero}.  The number \(S0\) is represented by \texttt{[1]}, and \(S(S0)\)
by \texttt{[[1]]}.  Raw scalar \(1\) in a packet's mode field is instead
an instruction marker; the field's role distinguishes the two uses.

Let \(R\) map an old canonical expression to the new one by replacing
each numerical-zero leaf \(0\) by \(1\), leaving every constructor,
protected coordinate, and rigid parameter unchanged.  This operation
is applied to expression fields.  A raw zero used as an instruction
marker remains zero.

\begin{lemma}[Uniform truthiness]
\label{truthy:positive}
Every canonical numerical expression has a truthy raw representation.
Every canonical formula is a nonempty array.  The old and new
representations decode to the same mathematical expression.
\end{lemma}

\begin{proof}
Canonical zero is the nonzero scalar \(1\).  A rigid parameter is a
strictly negative integer; a protected-coordinate or substitution-hole
occurrence has index at least \(3\).  Successor and iteration are
nonempty arrays.  These are all numerical constructors.
Equality, \(\Q\), and \(\D\) are respectively binary, ternary, and
unary nonempty arrays.  Finally, coordinates \(0\) and \(1\) already
have the same zero meaning in every numerical environment, and \(R\)
changes neither a binder nor its environment.  Structural induction
gives the last assertion.
\end{proof}

Truthiness is a property of this finite syntax.  A truthy term can
denote zero, and a truthy equation can be false.

\subsection{Opening preserves the convention}

Write \(i_0\) for the inherited opener.  The new opener \(i_+\) differs
only in the branch for a scalar already admitted by the current floor
and depth:
\[
             \texttt{a}\quad\longmapsto\quad\texttt{a or 1}.
\]
The list traversal, protected-coordinate test, numerical image check,
and depth increments are unchanged.

\begin{lemma}[Canonical opening commutation]
\label{truthy:opening}
For a well-scoped canonical expression \(A\), at a meaningful sorted
root of depth \(d\), and a canonical numerical image \(t\),
\[
 i_+\bigl(R(A),R(t),d\bigr)
       =R\bigl(i_0(A,t,d)\bigr).
 \tag{Opening commutation}\label{truthy:commute}
\]
The no-image version holds whenever opening does not use an image.
Opening a canonical template with raw image \(0\) in the new codec
has the same result as opening it with raw image \(1\).
\end{lemma}

\begin{proof}
First prove the no-image statement for closed images by structural induction.
Then proceed by structural induction on the template. At a meaningful numerical position
the depth is odd and at least \(3\).  If the old leaf is zero, it is
protected; the new representation and the changed preserved-scalar
branch both give \(1\).  Every other protected leaf is unchanged.
An exposed leaf invokes the numerical-root check of the image, and
the already proved no-image statement applies to that check.

For a successor, the same two dummy zero slots are introduced on both
sides.  For an iterator, the same single slot is introduced, with its
inherited zero or state interpretation in the respective child.
Equality changes from formula to numerical depth; the two logical
binders add the same protected slots.  The recursive results therefore
assemble into exactly the same raw tree on both sides of
\eqref{truthy:commute}.  A canonical formula has no scalar formula
root, so scalar formula padding does not arise in this exact identity.

An image used in a numerical position is itself checked at depth \(3\).
At every admissible floor,
\[
                 i_+(0,\mathrm{None},3)=1
                   =i_+(1,\mathrm{None},3).
\]
The two raw images therefore have identical inserted results.  If an
image is unused, its spelling is irrelevant.
\end{proof}

\begin{lemma}[Semantic opening on accepted aliases]
\label{truthy:semantic}
Every successful new opening at a meaningful sorted root has the same
capture-free substitution interpretation as the inherited opener.
\end{lemma}

\begin{proof}
The only changed scalar result is a falsy value retained by the scope
test, such as \(0\), signed floating zero, or \texttt{false}.
At a meaningful numerical position that value denotes zero; so does
the returned coordinate \(1\).  At a meaningful scalar formula
position both denote truth.  A call below a constant-meaning surplus
node concerns an ignored annotation.  All other successful scalar
cases and all constructor clauses are unchanged.  The inherited
structural semantic-opening proof consequently applies with this
additional scalar case.
\end{proof}

The exact identity in Lemma~\ref{truthy:opening} is deliberately
restricted to canonical expressions.  For arbitrary administrative
arrays, a subsequent traversal can treat an ignored coordinate
differently.  Scalar formula padding also need not be a fixed point
of raw traversal.  Lemma~\ref{truthy:semantic}, rather than universal
raw idempotence, is the invariant used for these aliases.

Unlike reduction of \(I(0,a;v.h)\) to its seed, this recoding deletes
no syntax node or protected slot.  In particular, it needs no
de Bruijn downshift beneath nested iterators.

\subsection{Constructing the elimination major}

An instruction has four fields \([m,s,u,t]\).  When \(m\) is truthy,
construct
\[
 M(m,s,u,t)=
 \begin{cases}
 [s,u],&s\text{ truthy},\\
 [t,t,t],&s\text{ falsy and }t\text{ truthy},\\
 m,&s,t\text{ both falsy}.
 \end{cases}
 \tag{Constructed major}\label{truthy:major}
\]
The checker proves \(M\) recursively in the retained context before
using it.  If opening \(t\) at \(u\) matches the current goal, its
candidate replacement is the opening of \(t\) at \(s\).  The usual
common continuation exports three copies of the checked major.
The source-level Boolean spelling may retain the original goal when
the replacement is falsy; this omission of a conversion is harmless.

\begin{proposition}[Soundness of constructed majors]
\label{truthy:major-sound}
Replacing the old explicit-major pattern comparison by
\eqref{truthy:major} preserves soundness of the elimination branch.
\end{proposition}

\begin{proof}
If \(s\) is truthy, the checked major is exactly the equation \(s=u\).
A successful target comparison identifies the old goal with \(F(u)\),
where \(F\) is the decoded template \(t\).
The continuation's changed goal is \(F(s)\), so ordinary equality
transport proves \(F(u)\).  There is no additional assertion that
some separately supplied formula has these endpoints.

If \(s\) is falsy and \(t\) truthy, the checked major is exactly
\(\Q_yF(y)\).  A successfully used falsy numerical image denotes
zero.  Thus the proposed minor is \(F(0)\), and the target is \(F(u)\);
this is the existing closure-elimination rule.  An unsupported falsy
JSON value rejects when a relevant numerical occurrence uses it.
When the image is unused, zero may be chosen as its representative.

Suppose finally that \(s,t\) are falsy.  The supplied major is checked,
and the existing triple export gives paired elimination when it is
unary.  Other nonempty array majors export a surplus truth formula.
There is no extra substantive target-conversion rule in this case:
a successfully opened falsy formula template denotes truth, so a
matching target already denotes truth.  For canonical nonempty goals,
the ordinary choice \(t=0\) yields scalar formula padding and cannot
match the goal at all.  Unsupported falsy values either reject or are
irrelevant.  This exhausts the branches.
\end{proof}

The freshness argument is unchanged: retained goals and hypotheses
are checked before a new local parameter is admitted.  Constructing
a major does not make it a hypothesis; its recursive proof remains
necessary.  Local allocation versus a global allocation counter is
a separate implementation choice, justified by the inherited scope
argument.

\subsection{Completeness without zero repairs}

\begin{theorem}[Direct canonical compilation]
\label{truthy:complete}
Every native derivation using the unrestricted I0 and IS computation
schemas has a canonical certificate under the new representation.
No computation node requires a proved nonempty replacement for zero.
On decoded terms and formulas, this compilation is the identity.
\end{theorem}

\begin{proof}
Use \(R\) on expression fields and preserve instruction markers.
The following packets implement the primitive inferences; fields
without a mathematical role are fixed as displayed:
\[
\begin{array}{l|l}
\text{inference}&(m,s,u,t)\\ \hline
\text{transport through }F&(1,R(s),R(u),R(F))\\
\Q\text{-elimination from }\Q B&(1,0,R(u),R(B))\\
\D\text{-elimination}&(R(\D B),0,0,0)\\
\mathrm{I0}\text{ with seed }a\text{ and update }h
 &(0,0,R(a),R(h))\\
\mathrm{IS}\text{ at predecessor }n
 &(0,R(n),R(a),R(h))\\
\mathrm{DI\text{-}case}\text{ with witness }w
 &([],0,0,R(w))\\
\mathrm{DI\text{-}progress}\text{ with witness }w
 &(0,0,0,R(w)).
\end{array}
\]
The paired matrix and its variable environments are those specified
by the goal, as in the inherited certificate compiler.

By Lemma~\ref{truthy:positive}, an equality source \(R(s)\) is always
truthy, including when \(s=0\).  A canonical closure matrix is
nonempty.  The constructed majors therefore select exactly the
required equality and closure theorems.  Their proof premises are
compiled recursively; Lemma~\ref{truthy:opening} gives the required
minor conclusions.  For paired elimination, the checked unary major
exports exactly its witness closure.

The initializer \(R(a)\) is always truthy, so every computation is
distinguished from paired introduction.  In I0, the generated raw
zero counter is converted to the canonical scalar \(1\) when the
supplied equation is opened.  In IS, even predecessor zero has
truthy representation \(1\), so the generated counter is
\texttt{[1]}, precisely the canonical successor of zero.
The update application commutes with \(R\) by
Lemma~\ref{truthy:opening}.  These are direct instances of the two
computation schemas.

The two introduction modes retain distinct falsy markers.
Their zero-instance images become canonical zero automatically;
fresh negative parameter images remain unchanged.  Recoding leaves
the matrix's root constructor unchanged, so progress remains
restricted to literal equations.  Hypothesis closure follows from
the same opening-commutation lemma.  Induction on the native
derivation now supplies the complete certificate.  Extra surplus
truth entries exported by equality or closure majors need not be
used by that certificate.

For noncanonical computation packets, the inherited known-zero-instance
argument still supplies soundness: specialize all newly admitted
parameters to numerical zero and cut in the resulting computation
theorem.  Lemma~\ref{truthy:semantic} preserves that argument.  It
does not claim that a generated hypothesis is uniformly true at a
new arbitrary parameter.

Finally, Lemma~\ref{truthy:positive} says that decoding every translated
expression returns the original expression.  The compilation therefore
changes neither its numerical interpretation nor its formula
interpretation.
\end{proof}

\begin{corollary}[Transfer of the full interpretation]
The canonical consequence relation checked by the new presentation
is the inherited native consequence relation.  Its whole-domain
bi-interpretation and equiconsistency with \(\ipra\) transfer unchanged,
including the formula return interpretations.
\end{corollary}

\begin{proof}
Theorem~\ref{truthy:complete} proves completeness on the canonical
image.  Lemma~\ref{truthy:semantic},
Proposition~\ref{truthy:major-sound}, and the inherited arithmetic,
introduction, and scope arguments prove soundness.
The decoding map is the identity on the native expressions, so
compose it with the existing two interpretations.  Their internal
proofs of both composites remain the same.
\end{proof}

\subsection{Exact certificate-size benefits}

Let \(|X|\) denote the UTF-8 byte length of compact JSON with no
whitespace, and use one-character scalar flags consistently in the
compared packets.  Then
\[
 |[X,Y]|=|X|+|Y|+3,\qquad
 |[B,B,B]|=3|B|+4.
\]
Replacing an equality major field \([s,u]\) by flag \(1\), while retaining
the endpoint fields already present in the packet, saves exactly
\[
                  \boxed{|s|+|u|+2}
 \tag{Equality saving}\label{truthy:eq-saving}
\]
bytes.  Replacing a closure major field \([B,B,B]\) by flag \(1\),
while retaining its matrix field, saves exactly
\[
                       \boxed{3|B|+3}
 \tag{Closure saving}\label{truthy:q-saving}
\]
bytes.  No inference node or proof premise is removed by these two
field changes.  The paired-major field is unchanged.

For the same native proof tree with the same parameter allocation,
zero recoding changes a one-character digit to another one-character
digit.  Thus the two displayed savings add exactly over its equality
and closure elimination packets.  The outer certificate-array
punctuation is unchanged because the packet count is unchanged.
This exact comparison can, for example, be applied to a common proof
already satisfying the old initializer restriction.  It also applies
to the fully repaired old proof before removing its repairs.

There is a separate proof-node benefit.  In the certified-zero
construction let
\[
                      C=I(S(S0),S0;v.0).
\]
The supplied repair replaces an excluded computation leaf by a
transport with the two computation premises \(C=0\) and \(E(C)\).
It therefore replaces one node by three.  If an unrestricted
derivation has \(N\) non-hypothesis nodes and \(r\) such computation
leaves, the displayed tree transformation has
\[
          N_{\mathrm{old\ repair}}=N+2r,
          \qquad N_{\mathrm{new}}=N,\qquad N+2r\leq 3N.
 \tag{Tree counts}\label{truthy:tree-count}
\]
The new construction directly serializes the original computation
nodes.  Where neither compiler closes an additional goal by a
previously retained hypothesis, the same equations count instruction
packets.  Context reuse can remove packets from either serialization,
so \eqref{truthy:tree-count} is not a theorem about shortest accepted
certificates.

If one shares a single \(C=0\) node in a DAG and \(r>0\), the old
displayed repair instead has \(N+r+1\) distinct nodes.  The new proof
still uses \(N\).  This is again a comparison of the specified
constructions, not a lower bound on all possible old derivations.

Both transformations have linear node overhead.  Constructing majors
copies already present fields at most three times, and \(C\) is a
fixed finite term.  Their syntactic work is consequently linear in
the fully expanded encoded proof data, apart from the ordinary
integer-name bookkeeping.  This statement does not remove the
inherited triplicate-closure codec's potential exponential expansion
relative to a different, shared representation of formulas.
When removal of repairs changes numerical parameter labels, exact
whole-file byte comparisons must also account for that renaming.

\paragraph{A complete small example.}
For \(I(0,0;v.0)=0\), the following compact old certificate has three
instruction packets and \(97\) bytes.  It uses the supplied
three-node repair and a scalar-zero computation marker accepted by
the old checker:
\begin{verbatim}
[[0,0,[[[0]],[0],0],0],[0,[0],[0],0],
 [[[[[0]],[0],0],0],[[[0]],[0],0],0,[[0,5,0],3]],
 [[0,0,0],0]]
\end{verbatim}
The line breaks above are for display and are excluded from the
measurement.  The new certificate is
\begin{verbatim}
[[0,0,1,1],[[1,1,1],1]]
\end{verbatim}
with one instruction packet and \(23\) bytes.  Both complete programs
accepted their respective certificate.  The mathematical conclusion
is identical under the two decoders.  The \(74\)-byte reduction is
exact for this pair; it does not assert that \(97\) bytes is the
minimum old certificate.

\subsection{What the result establishes}

The corpus already recognizes nonempty zero terms and scalar zero
aliases.  The new combination is the uniform canonical choice of a
truthy scalar zero, its preservation when the checker creates base
instances, and its use in constructing elimination majors.  The
proved benefits are direct zero computation and removal of redundant
major fields.  These are changes to proof presentation.

The native axiomatic basis is unchanged.  Kernel-size measurements
must therefore distinguish this contribution from the inherited
scope discipline, Boolean source factoring, and other implementation
choices.  Neither the representation theorem nor the certificate
counts establish an absolute simplest axiomatization or a major new
foundation of arithmetic.

